% Options for packages loaded elsewhere
\PassOptionsToPackage{unicode}{hyperref}
\PassOptionsToPackage{hyphens}{url}
\PassOptionsToPackage{dvipsnames,svgnames*,x11names*}{xcolor}
%
\documentclass[
  10pt,
]{article}
\usepackage{lmodern}
\usepackage{amssymb,amsmath}
\usepackage{ifxetex,ifluatex}
\ifnum 0\ifxetex 1\fi\ifluatex 1\fi=0 % if pdftex
  \usepackage[T1]{fontenc}
  \usepackage[utf8]{inputenc}
  \usepackage{textcomp} % provide euro and other symbols
\else % if luatex or xetex
  \usepackage{unicode-math}
  \defaultfontfeatures{Scale=MatchLowercase}
  \defaultfontfeatures[\rmfamily]{Ligatures=TeX,Scale=1}
  \setmainfont[]{DejaVu Serif}
  \setmonofont[]{DejaVu Sans Mono}
\fi
% Use upquote if available, for straight quotes in verbatim environments
\IfFileExists{upquote.sty}{\usepackage{upquote}}{}
\IfFileExists{microtype.sty}{% use microtype if available
  \usepackage[]{microtype}
  \UseMicrotypeSet[protrusion]{basicmath} % disable protrusion for tt fonts
}{}
\makeatletter
\@ifundefined{KOMAClassName}{% if non-KOMA class
  \IfFileExists{parskip.sty}{%
    \usepackage{parskip}
  }{% else
    \setlength{\parindent}{0pt}
    \setlength{\parskip}{6pt plus 2pt minus 1pt}}
}{% if KOMA class
  \KOMAoptions{parskip=half}}
\makeatother
\usepackage{xcolor}
\IfFileExists{xurl.sty}{\usepackage{xurl}}{} % add URL line breaks if available
\IfFileExists{bookmark.sty}{\usepackage{bookmark}}{\usepackage{hyperref}}
\hypersetup{
  colorlinks=true,
  linkcolor=Maroon,
  filecolor=Maroon,
  citecolor=Blue,
  urlcolor=Blue,
  pdfcreator={LaTeX via pandoc}}
\urlstyle{same} % disable monospaced font for URLs
\usepackage[margin=0.6in,landscape]{geometry}
\usepackage{longtable,booktabs}
% Correct order of tables after \paragraph or \subparagraph
\usepackage{etoolbox}
\makeatletter
\patchcmd\longtable{\par}{\if@noskipsec\mbox{}\fi\par}{}{}
\makeatother
% Allow footnotes in longtable head/foot
\IfFileExists{footnotehyper.sty}{\usepackage{footnotehyper}}{\usepackage{footnote}}
\makesavenoteenv{longtable}
\setlength{\emergencystretch}{3em} % prevent overfull lines
\providecommand{\tightlist}{%
  \setlength{\itemsep}{0pt}\setlength{\parskip}{0pt}}
\setcounter{secnumdepth}{-\maxdimen} % remove section numbering
\renewcommand{\arraystretch}{1.6}

\author{}
\date{}

\begin{document}

\hypertarget{phase-0.1-triglog-table-interpolation}{%
\section{Phase 0.1 --- Trig/Log Table
Interpolation}\label{phase-0.1-triglog-table-interpolation}}

\textbf{Type:} Tool drill, no mission narrative. \textbf{Target time:}
\textasciitilde30 min if interpolation is already comfortable; longer if
not. \textbf{Prerequisite:} 0.0 Slide rule basics. \textbf{Assumes:} You
have a 4- or 5-place log table and a matching trig table (sine/tangent,
tenths of a degree). The excerpts below are representative of real
period tables in layout and precision --- treat them as if torn from an
actual reference; a period-correct edition (CRC Standard Mathematical
Tables or similar) will be linked in \texttt{reference/tables.md} per
stage as sourcing is done.

\hypertarget{why-this-matters}{%
\subsection{Why this matters}\label{why-this-matters}}

A table only gives you exact values at its tabulated points. Everything
between two entries has to be interpolated --- and every
hand-computation procedure from this era (Gauss's method,
Kepler's-equation solving, orbit determination worksheets) leans on this
constantly. Get sloppy here and the error propagates into every later
step of a multi-stage hand solution.

\hypertarget{method-linear-interpolation}{%
\subsection{Method: linear
interpolation}\label{method-linear-interpolation}}

Given a table with entries at x₀ and x₁ = x₀ + Δx, and a target value x
between them:

\begin{verbatim}
f(x) ≈ f(x₀) + [(x - x₀) / Δx] · [f(x₁) - f(x₀)]
\end{verbatim}

In words: find how far across the interval your target sits (as a
fraction), then move that same fraction of the way between the two
tabulated values. This works well when the function is close to a
straight line over the interval --- which is the whole design goal of a
good table (fine enough spacing that linear interpolation is
trustworthy). It breaks down where the function curves sharply within
one interval --- see Drill 3 below.

\hypertarget{drill-1-log-table-interpolation}{%
\subsection{Drill 1 --- Log table
interpolation}\label{drill-1-log-table-interpolation}}

Table excerpt: log₁₀(x) for x = 1.00 to 1.10

\begin{longtable}[]{@{}ll@{}}
\toprule
x & log₁₀(x)\tabularnewline
\midrule
\endhead
1.00 & 0.00000\tabularnewline
1.01 & 0.00432\tabularnewline
1.02 & 0.00860\tabularnewline
1.03 & 0.01284\tabularnewline
1.04 & 0.01703\tabularnewline
1.05 & 0.02119\tabularnewline
1.06 & 0.02531\tabularnewline
1.07 & 0.02938\tabularnewline
1.08 & 0.03342\tabularnewline
1.09 & 0.03743\tabularnewline
1.10 & 0.04139\tabularnewline
\bottomrule
\end{longtable}

Interpolate:

\begin{enumerate}
\def\labelenumi{\arabic{enumi}.}
\tightlist
\item
  log₁₀(1.034)
\item
  log₁₀(1.067)
\item
  log₁₀(1.023)
\item
  log₁₀(1.089)
\end{enumerate}

\hypertarget{drill-2-trig-table-interpolation}{%
\subsection{Drill 2 --- Trig table
interpolation}\label{drill-2-trig-table-interpolation}}

Table excerpt: sin(θ) for θ = 23.0° to 24.0°

\begin{longtable}[]{@{}ll@{}}
\toprule
θ & sin(θ)\tabularnewline
\midrule
\endhead
23.0° & 0.39073\tabularnewline
23.1° & 0.39234\tabularnewline
23.2° & 0.39394\tabularnewline
23.3° & 0.39555\tabularnewline
23.4° & 0.39715\tabularnewline
23.5° & 0.39875\tabularnewline
23.6° & 0.40035\tabularnewline
23.7° & 0.40195\tabularnewline
23.8° & 0.40354\tabularnewline
23.9° & 0.40514\tabularnewline
24.0° & 0.40674\tabularnewline
\bottomrule
\end{longtable}

Interpolate:

\begin{enumerate}
\def\labelenumi{\arabic{enumi}.}
\tightlist
\item
  sin(23.34°)
\item
  sin(23.78°)
\item
  sin(23.06°)
\end{enumerate}

\hypertarget{drill-3-when-linear-interpolation-isnt-good-enough}{%
\subsection{Drill 3 --- When linear interpolation isn't good
enough}\label{drill-3-when-linear-interpolation-isnt-good-enough}}

Table excerpt: tan(θ) near a steep region

\begin{longtable}[]{@{}ll@{}}
\toprule
θ & tan(θ)\tabularnewline
\midrule
\endhead
80.0° & 5.67128\tabularnewline
80.5° & 5.97576\tabularnewline
81.0° & 6.31375\tabularnewline
\bottomrule
\end{longtable}

Notice the differences aren't constant (0.30448, then 0.33799) --- the
function is curving noticeably even over half-degree steps. This is what
``linear interpolation isn't accurate enough'' looks like in practice.

\begin{enumerate}
\def\labelenumi{\arabic{enumi}.}
\tightlist
\item
  Interpolate tan(80.3°) two ways: (a) using only the 80.0° and 81.0°
  entries (coarse table, 1° spacing), and (b) using the 80.0° and 80.5°
  entries (finer table, 0.5° spacing). Record both.
\item
  Compare the two answers. The gap between them is the real cost of
  using too coarse a table near a steep region of the function --- this
  is exactly why period trig tables near 80--90° were often tabulated at
  finer spacing (arcminutes, not tenths of a degree), or came with a
  second-difference correction column for exactly this situation.
\end{enumerate}

\end{document}
